\section{Introduction}

SAM~\cite{kirillov2023sam} was trained on over one billion masks and segments almost anything---except when the object has evolved not to be seen. On COD10K's camouflaged test set~\cite{fan2020cod10k}, SAM ViT-H scores 57.2\% mIoU with center-of-mass prompts and 20.5\% with edge prompts (\cref{tab:main}). Camouflaged organisms (insects mimicking bark, fish blending into coral) defeat visual detection by minimizing texture contrast with their surroundings, and SAM's decoder has never been trained to read these signals as object boundaries.

We ask: \textbf{is the decoder the bottleneck?} The ViT-H encoder, pre-trained on 11M images, may already encode texture gradients that separate camouflaged foreground from background. The decoder, trained on general segmentation, may simply lack the mapping from those gradients to object boundaries. If so, fine-tuning the 4M-parameter mask decoder alone should be enough.

It is. Decoder-only fine-tuning on 3,040 COD10K training images raises mIoU from 57.2\% to 73.8\%. Following SAM's standard evaluation protocol~\cite{kirillov2023sam}, prompts are derived from ground-truth masks to simulate an idealized user click. The improvement is \textbf{prompt-robust}: it transfers across center, edge, grid, and random prompt strategies (\cref{tab:improvement}), with edge prompts showing the largest gain: 20.5\% to 69.4\% (+238\% relative). This robustness is not accidental. Ablating the training prompt mix (\cref{sec:ablation_prompt}) shows that mixed point/box training outperforms point-only by up to 4 mIoU and box-only by 19--37 mIoU.

We also introduce a \textbf{Learnable Local Preprocessing Module (LLPM)}: a 31K-parameter module placed \emph{before} the frozen encoder that amplifies boundary signals in input space. The LLPM adds +2.6 mIoU on edge prompts over the decoder-only baseline, targeting the boundary-ambiguous failures identified by our error analysis (\cref{sec:error_analysis}).

Analyzing the remaining failures shows they are not random: dark images, tiny objects ($<$1\% foreground pixels), and thin/elongated shapes account for most low-IoU predictions.

\noindent\textbf{Contributions:}
\begin{enumerate}
    \item Decoder-only adaptation for promptable COD: fine-tuning SAM's mask decoder alone yields +16.6 mIoU on 2,026 test images with no architecture changes, 90-minute training, and a 16MB checkpoint.
    \item Prompt robustness analysis: mixed-prompt training reduces sensitivity to prompt type from 3.5$\times$ to 1.1$\times$, validated by point-only and box-only ablations.
    \item A 31K-parameter LLPM that amplifies boundary signals before the frozen encoder, adding +2.6 mIoU on edge prompts.
    \item Error analysis of failure modes (dark, tiny, thin, disjoint) that links LLPM gains to specific failure categories.
\end{enumerate}
